\chapter{Unit 2: Natural Resources and Ecosystem Dynamics}

\section{Section A: Multiple Choice Practice Questions (60 Curated Problems)}

\noindent\textbf{Q1. What are natural resources?} \hfill [Tag: Introduction to natural resources | Easy]
\begin{enumerate}[label=(\alph*),topsep=2pt,itemsep=1pt,partopsep=0pt,parsep=0pt]
    \item {Goods imported from other countries}
    \item {Products made only in factories}
    \item {Materials obtained from nature}
    \item {Objects created by computer programs}
\end{enumerate}
\vspace{2mm}

\noindent\textbf{Q2. Which of the following is a renewable resource?} \hfill [Tag: Renewable and non-renewable resources | Easy]
\begin{enumerate}[label=(\alph*),topsep=2pt,itemsep=1pt,partopsep=0pt,parsep=0pt]
    \item {Solar energy}
    \item {Coal}
    \item {Petroleum}
    \item {Natural gas}
\end{enumerate}
\vspace{2mm}

\noindent\textbf{Q3. Which resource is classified as non-renewable?} \hfill [Tag: Renewable and non-renewable resources | Easy]
\begin{enumerate}[label=(\alph*),topsep=2pt,itemsep=1pt,partopsep=0pt,parsep=0pt]
    \item {Rainwater}
    \item {Wind}
    \item {Sunlight}
    \item {Coal}
\end{enumerate}
\vspace{2mm}

\noindent\textbf{Q4. Which resource provides habitat for many wild plants and animals?} \hfill [Tag: Land, water, forest and energy resources | Easy]
\begin{enumerate}[label=(\alph*),topsep=2pt,itemsep=1pt,partopsep=0pt,parsep=0pt]
    \item {Petroleum}
    \item {Mineral ore}
    \item {Natural gas}
    \item {Forest}
\end{enumerate}
\vspace{2mm}

\noindent\textbf{Q5. Which activity uses the largest share of freshwater in many countries?} \hfill [Tag: Land, water, forest and energy resources | Easy]
\begin{enumerate}[label=(\alph*),topsep=2pt,itemsep=1pt,partopsep=0pt,parsep=0pt]
    \item {Agriculture}
    \item {Communication}
    \item {Transportation}
    \item {Construction}
\end{enumerate}
\vspace{2mm}

\noindent\textbf{Q6. Which energy resource is obtained from moving air?} \hfill [Tag: Land, water, forest and energy resources | Easy]
\begin{enumerate}[label=(\alph*),topsep=2pt,itemsep=1pt,partopsep=0pt,parsep=0pt]
    \item {Tidal energy}
    \item {Geothermal energy}
    \item {Wind energy}
    \item {Coal energy}
\end{enumerate}
\vspace{2mm}

\noindent\textbf{Q7. Which method helps reduce soil erosion on steep slopes?} \hfill [Tag: Associated problems and remedial measures of natural resources | Easy]
\begin{enumerate}[label=(\alph*),topsep=2pt,itemsep=1pt,partopsep=0pt,parsep=0pt]
    \item {Overgrazing}
    \item {Deforestation}
    \item {Terrace farming}
    \item {Open mining}
\end{enumerate}
\vspace{2mm}

\noindent\textbf{Q8. What is a major result of excessive removal of trees?} \hfill [Tag: Associated problems and remedial measures of natural resources | Easy]
\begin{enumerate}[label=(\alph*),topsep=2pt,itemsep=1pt,partopsep=0pt,parsep=0pt]
    \item {Afforestation}
    \item {Water harvesting}
    \item {Deforestation}
    \item {Soil formation}
\end{enumerate}
\vspace{2mm}

\noindent\textbf{Q9. Which practice helps replenish groundwater?} \hfill [Tag: Associated problems and remedial measures of natural resources | Easy]
\begin{enumerate}[label=(\alph*),topsep=2pt,itemsep=1pt,partopsep=0pt,parsep=0pt]
    \item {Rainwater harvesting}
    \item {Clearing vegetation}
    \item {Burning fossil fuels}
    \item {Paving open land}
\end{enumerate}
\vspace{2mm}

\noindent\textbf{Q10. Which individual action directly conserves water?} \hfill [Tag: Role of individual in conservation of natural resources | Easy]
\begin{enumerate}[label=(\alph*),topsep=2pt,itemsep=1pt,partopsep=0pt,parsep=0pt]
    \item {Using disposable plates}
    \item {Burning household waste}
    \item {Repairing leaking taps}
    \item {Leaving lights switched on}
\end{enumerate}
\vspace{2mm}

\noindent\textbf{Q11. Which action helps conserve energy at home?} \hfill [Tag: Role of individual in conservation of natural resources | Easy]
\begin{enumerate}[label=(\alph*),topsep=2pt,itemsep=1pt,partopsep=0pt,parsep=0pt]
    \item {Keeping doors open while cooling}
    \item {Running empty washing machines}
    \item {Switching off unused lights}
    \item {Using appliances throughout the day}
\end{enumerate}
\vspace{2mm}

\noindent\textbf{Q12. Which of the following is a biotic component of an ecosystem?} \hfill [Tag: Structure and function of ecosystem | Easy]
\begin{enumerate}[label=(\alph*),topsep=2pt,itemsep=1pt,partopsep=0pt,parsep=0pt]
    \item {Sunlight}
    \item {Water}
    \item {Soil}
    \item {Plant}
\end{enumerate}
\vspace{2mm}

\noindent\textbf{Q13. What is the main role of decomposers in an ecosystem?} \hfill [Tag: Structure and function of ecosystem | Easy]
\begin{enumerate}[label=(\alph*),topsep=2pt,itemsep=1pt,partopsep=0pt,parsep=0pt]
    \item {Breaking down dead matter}
    \item {Hunting living herbivores}
    \item {Preventing rainfall from occurring}
    \item {Producing energy from sunlight}
\end{enumerate}
\vspace{2mm}

\noindent\textbf{Q14. Which of the following is an aquatic ecosystem?} \hfill [Tag: Types of ecosystem | Easy]
\begin{enumerate}[label=(\alph*),topsep=2pt,itemsep=1pt,partopsep=0pt,parsep=0pt]
    \item {Grassland}
    \item {Pond}
    \item {Forest}
    \item {Desert}
\end{enumerate}
\vspace{2mm}

\noindent\textbf{Q15. Which ecosystem usually receives very little rainfall?} \hfill [Tag: Types of ecosystem | Easy]
\begin{enumerate}[label=(\alph*),topsep=2pt,itemsep=1pt,partopsep=0pt,parsep=0pt]
    \item {Rainforest}
    \item {Wetland}
    \item {Estuary}
    \item {Desert}
\end{enumerate}
\vspace{2mm}

\noindent\textbf{Q16. What is the primary source of energy for most ecosystems?} \hfill [Tag: Energy flow in an ecosystem | Easy]
\begin{enumerate}[label=(\alph*),topsep=2pt,itemsep=1pt,partopsep=0pt,parsep=0pt]
    \item {The Sun}
    \item {The wind}
    \item {The Moon}
    \item {The soil}
\end{enumerate}
\vspace{2mm}

\noindent\textbf{Q17. In the food chain grass $\to$ deer $\to$ tiger , which organism is the producer?} \hfill [Tag: Food chains and food webs | Easy]
\begin{enumerate}[label=(\alph*),topsep=2pt,itemsep=1pt,partopsep=0pt,parsep=0pt]
    \item {Tiger}
    \item {Grass}
    \item {None}
    \item {Deer}
\end{enumerate}
\vspace{2mm}

\noindent\textbf{Q18. What does a food web represent?} \hfill [Tag: Food chains and food webs | Easy]
\begin{enumerate}[label=(\alph*),topsep=2pt,itemsep=1pt,partopsep=0pt,parsep=0pt]
    \item {Interconnected food chains}
    \item {Only decomposer activity}
    \item {Only plant relationships}
    \item {A single feeding pathway}
\end{enumerate}
\vspace{2mm}

\noindent\textbf{Q19. Which organisms usually form the base of an ecological pyramid?} \hfill [Tag: Ecological pyramids | Easy]
\begin{enumerate}[label=(\alph*),topsep=2pt,itemsep=1pt,partopsep=0pt,parsep=0pt]
    \item {Decomposers}
    \item {Herbivores}
    \item {Producers}
    \item {Carnivores}
\end{enumerate}
\vspace{2mm}

\noindent\textbf{Q20. What is ecological succession?} \hfill [Tag: Ecological succession | Easy]
\begin{enumerate}[label=(\alph*),topsep=2pt,itemsep=1pt,partopsep=0pt,parsep=0pt]
    \item {Sudden transfer of energy between organisms}
    \item {Daily movement of animals between habitats}
    \item {Seasonal change in local weather conditions}
    \item {Gradual change in a biological community}
\end{enumerate}
\vspace{2mm}

\noindent\textbf{Q21. A village depends on groundwater for drinking and irrigation. Which statement best explains why groundwater is considered a natural resource?} \hfill [Tag: Introduction to natural resources | Medium]
\begin{enumerate}[label=(\alph*),topsep=2pt,itemsep=1pt,partopsep=0pt,parsep=0pt]
    \item {It is produced only through industrial processes}
    \item {It is available in unlimited amounts everywhere}
    \item {It is obtained from nature and supports human needs}
    \item {It has no connection with ecological systems}
\end{enumerate}
\vspace{2mm}

\noindent\textbf{Q22. Which situation best illustrates that a renewable resource can still become scarce?} \hfill [Tag: Renewable and non-renewable resources | Medium]
\begin{enumerate}[label=(\alph*),topsep=2pt,itemsep=1pt,partopsep=0pt,parsep=0pt]
    \item {Solar radiation reaches Earth during the daytime}
    \item {Wind turbines generate electricity in a coastal area}
    \item {Coal is extracted from a newly discovered mine}
    \item {A forest is harvested faster than it regenerates}
\end{enumerate}
\vspace{2mm}

\noindent\textbf{Q23. Which agricultural practice is most effective in reducing soil erosion on a steep slope?} \hfill [Tag: Land, water, forest and energy resources | Medium]
\begin{enumerate}[label=(\alph*),topsep=2pt,itemsep=1pt,partopsep=0pt,parsep=0pt]
    \item {Deep ploughing down the slope}
    \item {Removing vegetation before monsoon}
    \item {Increasing irrigation during heavy rainfall}
    \item {Contour ploughing across the slope}
\end{enumerate}
\vspace{2mm}

\noindent\textbf{Q24. A town experiences falling groundwater levels despite receiving normal annual rainfall. Which change would most directly help restore groundwater?} \hfill [Tag: Land, water, forest and energy resources | Medium]
\begin{enumerate}[label=(\alph*),topsep=2pt,itemsep=1pt,partopsep=0pt,parsep=0pt]
    \item {Increasing the use of deep bore wells}
    \item {Diverting wastewater into nearby natural lakes}
    \item {Constructing rooftop rainwater harvesting systems}
    \item {Paving additional open spaces with concrete}
\end{enumerate}
\vspace{2mm}

\noindent\textbf{Q25. A reservoir is receiving large amounts of soil from surrounding hills, reducing its storage capacity. Which remedial measure is most appropriate?} \hfill [Tag: Associated problems and remedial measures of natural resources | Medium]
\begin{enumerate}[label=(\alph*),topsep=2pt,itemsep=1pt,partopsep=0pt,parsep=0pt]
    \item {Replacing forests with intensive crop fields}
    \item {Increasing the speed of water entering the reservoir}
    \item {Planting vegetation and building check dams}
    \item {Removing all grasses from the catchment}
\end{enumerate}
\vspace{2mm}

\noindent\textbf{Q26. Which combination most effectively reduces the environmental problems caused by excessive dependence on fossil fuels?} \hfill [Tag: Associated problems and remedial measures of natural resources | Medium]
\begin{enumerate}[label=(\alph*),topsep=2pt,itemsep=1pt,partopsep=0pt,parsep=0pt]
    \item {Energy efficiency and renewable energy use}
    \item {Higher electricity demand and reduced public transport}
    \item {Greater coal consumption and larger vehicles}
    \item {More petroleum extraction and open burning}
\end{enumerate}
\vspace{2mm}

\noindent\textbf{Q27. Which individual action would conserve both water and energy in a household?} \hfill [Tag: Role of individual in conservation of natural resources | Medium]
\begin{enumerate}[label=(\alph*),topsep=2pt,itemsep=1pt,partopsep=0pt,parsep=0pt]
    \item {Repairing leaking taps promptly}
    \item {Using hot water for every cleaning task}
    \item {Washing vehicles daily with a hose}
    \item {Keeping taps running while brushing}
\end{enumerate}
\vspace{2mm}

\noindent\textbf{Q28. A student wants to reduce the amount of waste sent to a landfill. Which choice follows the waste-management hierarchy most effectively?} \hfill [Tag: Role of individual in conservation of natural resources | Medium]
\begin{enumerate}[label=(\alph*),topsep=2pt,itemsep=1pt,partopsep=0pt,parsep=0pt]
    \item {Burn all household waste in an open area}
    \item {Send recyclable materials to the landfill with other waste}
    \item {Buy disposable products because they require no cleaning}
    \item {Refuse unnecessary packaging and reuse containers}
\end{enumerate}
\vspace{2mm}

\noindent\textbf{Q29. In an ecosystem, decomposers are essential mainly because they} \hfill [Tag: Structure and function of ecosystem | Medium]
\begin{enumerate}[label=(\alph*),topsep=2pt,itemsep=1pt,partopsep=0pt,parsep=0pt]
    \item {Convert all sunlight directly into heat}
    \item {Consume only inorganic substances for energy}
    \item {Return nutrients to the environment}
    \item {Prevent producers from carrying out photosynthesis}
\end{enumerate}
\vspace{2mm}

\noindent\textbf{Q30. Which example correctly matches an ecosystem component with its role?} \hfill [Tag: Structure and function of ecosystem | Medium]
\begin{enumerate}[label=(\alph*),topsep=2pt,itemsep=1pt,partopsep=0pt,parsep=0pt]
    \item {Producer: converts solar energy into chemical energy}
    \item {Abiotic factor: feeds directly on herbivores}
    \item {Consumer: manufactures minerals from sunlight}
    \item {Decomposer: prevents nutrient cycling in soil}
\end{enumerate}
\vspace{2mm}

\noindent\textbf{Q31. A pond contains algae, fish, aquatic plants, microbes, water, dissolved salts, and sunlight. Which classification best describes this system?} \hfill [Tag: Types of ecosystem | Medium]
\begin{enumerate}[label=(\alph*),topsep=2pt,itemsep=1pt,partopsep=0pt,parsep=0pt]
    \item {A natural freshwater ecosystem}
    \item {An artificial desert ecosystem}
    \item {A natural marine ecosystem}
    \item {An artificial terrestrial ecosystem}
\end{enumerate}
\vspace{2mm}

\noindent\textbf{Q32. Which feature most clearly distinguishes an artificial ecosystem from a natural ecosystem?} \hfill [Tag: Types of ecosystem | Medium]
\begin{enumerate}[label=(\alph*),topsep=2pt,itemsep=1pt,partopsep=0pt,parsep=0pt]
    \item {It lacks any flow of energy}
    \item {It has no interaction among organisms}
    \item {It requires regular human management}
    \item {It contains only living organisms}
\end{enumerate}
\vspace{2mm}

\noindent\textbf{Q33. If producers in an ecosystem store 10 , 000 units of energy, approximately how much energy may be available to primary consumers according to the ten-percent law?} \hfill [Tag: Energy flow in an ecosystem | Medium]
\begin{enumerate}[label=(\alph*),topsep=2pt,itemsep=1pt,partopsep=0pt,parsep=0pt]
    \item {9 ,}
    \item {1 ,}
    \item {100}
    \item {5 ,}
\end{enumerate}
\vspace{2mm}

\noindent\textbf{Q34. Why does energy flow through an ecosystem rather than cycle like nutrients?} \hfill [Tag: Energy flow in an ecosystem | Medium]
\begin{enumerate}[label=(\alph*),topsep=2pt,itemsep=1pt,partopsep=0pt,parsep=0pt]
    \item {Energy is lost as heat at each transfer}
    \item {Energy is created by consumers during feeding}
    \item {Energy returns completely to producers after use}
    \item {Energy is permanently stored in decomposers}
\end{enumerate}
\vspace{2mm}

\noindent\textbf{Q35. In the food chain grass $\to$ grasshopper $\to$ frog $\to$ snake, what is the likely immediate effect of a sharp decrease in frogs?} \hfill [Tag: Food chains and food webs | Medium]
\begin{enumerate}[label=(\alph*),topsep=2pt,itemsep=1pt,partopsep=0pt,parsep=0pt]
    \item {Grasshoppers decrease and grass increases}
    \item {Grass and grasshoppers both increase equally}
    \item {Snakes increase and grasshoppers decrease}
    \item {Grasshoppers increase and grass decreases}
\end{enumerate}
\vspace{2mm}

\noindent\textbf{Q36. Why is a food web generally more stable than a single food chain?} \hfill [Tag: Food chains and food webs | Medium]
\begin{enumerate}[label=(\alph*),topsep=2pt,itemsep=1pt,partopsep=0pt,parsep=0pt]
    \item {It eliminates competition among all organisms}
    \item {It prevents predators from feeding on consumers}
    \item {It provides alternative feeding pathways}
    \item {It transfers the same energy to every organism}
\end{enumerate}
\vspace{2mm}

\noindent\textbf{Q37. Which ecological pyramid is always upright in a normal ecosystem?} \hfill [Tag: Ecological pyramids | Medium]
\begin{enumerate}[label=(\alph*),topsep=2pt,itemsep=1pt,partopsep=0pt,parsep=0pt]
    \item {Pyramid of biomass in a pond}
    \item {Pyramid of biomass in a parasitic chain}
    \item {Pyramid of numbers in a forest}
    \item {Pyramid of energy}
\end{enumerate}
\vspace{2mm}

\noindent\textbf{Q38. An ecosystem has one large tree supporting many insects and birds. Which pyramid may have an inverted shape at the producer and primary-consumer levels?} \hfill [Tag: Ecological pyramids | Medium]
\begin{enumerate}[label=(\alph*),topsep=2pt,itemsep=1pt,partopsep=0pt,parsep=0pt]
    \item {Pyramid of productivity}
    \item {Pyramid of available sunlight}
    \item {Pyramid of numbers}
    \item {Pyramid of energy}
\end{enumerate}
\vspace{2mm}

\noindent\textbf{Q39. Which sequence best represents primary succession on newly exposed bare rock?} \hfill [Tag: Ecological succession | Medium]
\begin{enumerate}[label=(\alph*),topsep=2pt,itemsep=1pt,partopsep=0pt,parsep=0pt]
    \item {Mature forest, bare rock, grasses, and shrubs}
    \item {Herbivores, trees, lichens, and decomposers}
    \item {Lichens, grasses, shrubs, and trees}
    \item {Trees, shrubs, grasses, and lichens}
\end{enumerate}
\vspace{2mm}

\noindent\textbf{Q40. A grassland is destroyed by fire, but the soil and some seeds remain. The subsequent recovery is an example of} \hfill [Tag: Ecological succession | Medium]
\begin{enumerate}[label=(\alph*),topsep=2pt,itemsep=1pt,partopsep=0pt,parsep=0pt]
    \item {Primary succession}
    \item {Succession without pioneer species}
    \item {Secondary succession}
    \item {Hydrarch succession}
\end{enumerate}
\vspace{2mm}

\noindent\textbf{Q41. A coastal wetland supports fisheries, buffers storm surges, stores carbon, and provides habitat for migratory birds. A valuation study counts fish harvest revenue and avoided storm damage but excludes carbon storage and habitat support. Which conclusion is most defensible?} \hfill [Tag: Introduction to natural resources | Hard]
\begin{enumerate}[label=(\alph*),topsep=2pt,itemsep=1pt,partopsep=0pt,parsep=0pt]
    \item {The estimate overstates total economic value by combining direct and indirect uses}
    \item {The estimate represents non-use value because fisheries are harvested outside the wetland}
    \item {The estimate understates total economic value by omitting regulating and supporting services}
    \item {The estimate equals total economic value because only market benefits are measurable}
\end{enumerate}
\vspace{2mm}

\noindent\textbf{Q42. A confined aquifer receives negligible modern recharge, while a forest regrows naturally if annual harvest remains below net growth. Which classification best reflects resource dynamics rather than merely resource identity?} \hfill [Tag: Renewable and non-renewable resources | Hard]
\begin{enumerate}[label=(\alph*),topsep=2pt,itemsep=1pt,partopsep=0pt,parsep=0pt]
    \item {Both are renewable because each is part of a natural cycle}
    \item {The aquifer is effectively non-renewable, while the forest is conditionally renewable}
    \item {Both are non-renewable because human use reduces their standing stocks}
    \item {The aquifer is renewable, while the forest is non-renewable after any harvest}
\end{enumerate}
\vspace{2mm}

\noindent\textbf{Q43. A fish population follows logistic growth, G ( N ) = r N ( 1 − N / K ) . Ignoring age structure and environmental variability, at what population size is the theoretical maximum sustainable yield obtained?} \hfill [Tag: Renewable and non-renewable resources | Hard]
\begin{enumerate}[label=(\alph*),topsep=2pt,itemsep=1pt,partopsep=0pt,parsep=0pt]
    \item {N =}
    \item {N =}
    \item {N =}
    \item {N =}
\end{enumerate}
\vspace{2mm}

\noindent\textbf{Q44. An irrigated basin has rising groundwater, increasing soil salinity, and declining crop yields despite adequate rainfall. Which intervention most directly addresses the coupled cause of the degradation?} \hfill [Tag: Land, water, forest and energy resources | Hard]
\begin{enumerate}[label=(\alph*),topsep=2pt,itemsep=1pt,partopsep=0pt,parsep=0pt]
    \item {Replace annual crops with deeper-rooted crops without changing water inputs}
    \item {Install subsurface drainage and apply controlled leaching with suitable water}
    \item {Increase fertilizer application and retain the existing irrigation schedule}
    \item {Line field channels and discontinue all removal of saline drainage water}
\end{enumerate}
\vspace{2mm}

\noindent\textbf{Q45. A region must choose an electricity project using life-cycle analysis. Project P has low operating emissions but high construction emissions; Project Q has moderate operating emissions and low construction emissions. What information is essential for a valid climate comparison?} \hfill [Tag: Land, water, forest and energy resources | Hard]
\begin{enumerate}[label=(\alph*),topsep=2pt,itemsep=1pt,partopsep=0pt,parsep=0pt]
    \item {Only the rated capacity and maximum hourly output of each project}
    \item {Annual operating emissions without accounting for generated electricity}
    \item {Life-cycle emissions per unit of electricity over comparable service lifetimes}
    \item {Only the construction emissions and geographic area of each project}
\end{enumerate}
\vspace{2mm}

\noindent\textbf{Q46. Two forest plans remove the same timber volume. Plan I clear-cuts one contiguous block; Plan II removes scattered trees across a larger area while building many access roads. Which statement best explains why Plan II may produce greater ecological damage?} \hfill [Tag: Land, water, forest and energy resources | Hard]
\begin{enumerate}[label=(\alph*),topsep=2pt,itemsep=1pt,partopsep=0pt,parsep=0pt]
    \item {Road density can increase fragmentation, edge effects, erosion, and human access}
    \item {Scattered harvesting prevents invasive species from reaching forest interiors}
    \item {Selective removal always eliminates more canopy biomass than clear-cutting}
    \item {Clear-cut blocks cannot alter stream flow when their total area is limited}
\end{enumerate}
\vspace{2mm}

\noindent\textbf{Q47. A watershed experiences severe peak floods, reduced dry-season streamflow, and heavy sedimentation after conversion of upland forest to compacted cropland. Which restoration package best addresses all three symptoms?} \hfill [Tag: Associated problems and remedial measures of natural resources | Hard]
\begin{enumerate}[label=(\alph*),topsep=2pt,itemsep=1pt,partopsep=0pt,parsep=0pt]
    \item {Increase groundwater pumping, straighten channels, and burn crop residues}
    \item {Reforest steep slopes, restore riparian buffers, and reduce soil compaction}
    \item {Channelize streams, remove riparian vegetation, and deepen downstream reservoirs}
    \item {Construct higher levees, expand paved drainage, and dredge river bends}
\end{enumerate}
\vspace{2mm}

\noindent\textbf{Q48. A common pasture is degraded because each herder receives the full benefit of adding an animal but shares the resulting grazing cost with all users. Which policy most directly aligns individual incentives with the pasture's carrying capacity?} \hfill [Tag: Associated problems and remedial measures of natural resources | Hard]
\begin{enumerate}[label=(\alph*),topsep=2pt,itemsep=1pt,partopsep=0pt,parsep=0pt]
    \item {Establish enforceable community quotas with monitoring and graduated sanctions}
    \item {Replace rotational grazing with continuous grazing throughout the pasture}
    \item {Provide an unrestricted subsidy for every additional animal maintained}
    \item {Guarantee open access while periodically distributing supplementary fodder}
\end{enumerate}
\vspace{2mm}

\noindent\textbf{Q49. A household can either replace a functioning refrigerator immediately with a highly efficient model or continue using it for several years. Which criterion gives the most defensible conservation decision?} \hfill [Tag: Role of individual in conservation of natural resources | Hard]
\begin{enumerate}[label=(\alph*),topsep=2pt,itemsep=1pt,partopsep=0pt,parsep=0pt]
    \item {Retain the old model whenever manufacturing requires any non-renewable material}
    \item {Compare purchase prices while excluding electricity source and disposal impacts}
    \item {Choose replacement whenever the new model has a higher efficiency label}
    \item {Compare life-cycle energy and material impacts over equivalent years of service}
\end{enumerate}
\vspace{2mm}

\noindent\textbf{Q50. In a grassland experiment, decomposers are strongly suppressed while plant photosynthesis and herbivore abundance initially remain unchanged. Which effect is most likely after several growing seasons?} \hfill [Tag: Structure and function of ecosystem | Hard]
\begin{enumerate}[label=(\alph*),topsep=2pt,itemsep=1pt,partopsep=0pt,parsep=0pt]
    \item {Faster mineralization and a sustained rise in primary productivity}
    \item {Permanent accumulation of mineral nutrients in plant-available forms}
    \item {Immediate elimination of respiration from every trophic level}
    \item {Reduced nutrient availability and eventual decline in primary productivity}
\end{enumerate}
\vspace{2mm}

\noindent\textbf{Q51. An ecosystem has gross primary productivity of 2400 kJ m − 2 y r − 1 . Producers use 900 kJ m − 2 y r − 1 in respiration, and herbivores consume 300 kJ m − 2 y r − 1 of plant production. What are the net primary productivity and the plant production entering detrital or storage pathways?} \hfill [Tag: Structure and function of ecosystem | Hard]
\begin{enumerate}[label=(\alph*),topsep=2pt,itemsep=1pt,partopsep=0pt,parsep=0pt]
    \item {1500}
    \item {2400}
    \item {2100}
    \item {1500}
\end{enumerate}
\vspace{2mm}

\noindent\textbf{Q52. A shallow lake shifts from clear water with submerged plants to turbid water dominated by phytoplankton after nutrient enrichment. Nutrient inputs are then reduced slightly, but the turbid state persists. Which concept best explains this response?} \hfill [Tag: Types of ecosystem | Hard]
\begin{enumerate}[label=(\alph*),topsep=2pt,itemsep=1pt,partopsep=0pt,parsep=0pt]
    \item {Ecological equivalence caused by identical niches in both lake states}
    \item {Primary succession caused by complete removal of the lake's soil substrate}
    \item {Competitive release caused solely by the disappearance of all consumers}
    \item {Alternative stable states maintained by feedbacks and ecological hysteresis}
\end{enumerate}
\vspace{2mm}

\noindent\textbf{Q53. Compared with an oligotrophic lake of similar depth, a eutrophic lake undergoes strong summer stratification and receives high organic inputs. Which oxygen pattern is most likely before autumn mixing?} \hfill [Tag: Types of ecosystem | Hard]
\begin{enumerate}[label=(\alph*),topsep=2pt,itemsep=1pt,partopsep=0pt,parsep=0pt]
    \item {Low hypolimnetic oxygen because stratification limits replacement while decomposition consumes it}
    \item {Zero epilimnetic oxygen because photosynthesis is confined to deep water}
    \item {High hypolimnetic oxygen because decomposition decreases with organic input}
    \item {Uniform oxygen because thermal stratification accelerates vertical circulation}
\end{enumerate}
\vspace{2mm}

\noindent\textbf{Q54. Producers store 20 , 000 kJ as net production. Herbivores ingest 4 , 000 kJ , egest 1 , 200 kJ , and spend 2 , 240 kJ in respiration. What are the herbivores' assimilation efficiency and production efficiency?} \hfill [Tag: Energy flow in an ecosystem | Hard]
\begin{enumerate}[label=(\alph*),topsep=2pt,itemsep=1pt,partopsep=0pt,parsep=0pt]
    \item {20\%}
    \item {14\%}
    \item {70\%}
    \item {70\%}
\end{enumerate}
\vspace{2mm}

\noindent\textbf{Q55. A contaminant reduces producer respiration without changing gross photosynthesis, but it does not directly affect consumers. Which immediate ecosystem-level change is expected, assuming other losses remain constant?} \hfill [Tag: Energy flow in an ecosystem | Hard]
\begin{enumerate}[label=(\alph*),topsep=2pt,itemsep=1pt,partopsep=0pt,parsep=0pt]
    \item {Net primary productivity increases, potentially increasing energy available to consumers}
    \item {Consumer assimilation efficiency necessarily falls to the producer respiration rate}
    \item {Gross primary productivity increases by exactly the reduction in respiration}
    \item {Net primary productivity decreases because producer respiration is part of production}
\end{enumerate}
\vspace{2mm}

\noindent\textbf{Q56. In a rocky-shore food web, sea stars consume mussels, and mussels competitively exclude several other sessile species when abundant. Removing sea stars causes species richness to decline. Which mechanism best explains the result?} \hfill [Tag: Food chains and food webs | Hard]
\begin{enumerate}[label=(\alph*),topsep=2pt,itemsep=1pt,partopsep=0pt,parsep=0pt]
    \item {Predation suppresses a dominant competitor and permits species coexistence}
    \item {Mussels become primary producers after release from sea-star predation}
    \item {Top predators directly supply limiting nutrients to every sessile species}
    \item {Predator removal increases resource partitioning among all sessile consumers}
\end{enumerate}
\vspace{2mm}

\noindent\textbf{Q57. A pesticide enters a lake at low concentration, is persistent, lipid-soluble, and poorly metabolized. The food web is phytoplankton $\to$ zooplankton $\to$ small fish $\to$ fish-eating birds. Which observation would best demonstrate biomagnification rather than only bioaccumulation?} \hfill [Tag: Food chains and food webs | Hard]
\begin{enumerate}[label=(\alph*),topsep=2pt,itemsep=1pt,partopsep=0pt,parsep=0pt]
    \item {Phytoplankton contain pesticide even though the compound entered through runoff}
    \item {Lake sediment retains pesticide after concentrations decline in the overlying water}
    \item {Bird tissues contain higher concentrations than fish, which exceed zooplankton concentrations}
    \item {Individual zooplankton accumulate more pesticide as their exposure duration increases}
\end{enumerate}
\vspace{2mm}

\noindent\textbf{Q58. An open-ocean ecosystem has a smaller standing biomass of phytoplankton than of zooplankton at a given instant, yet phytoplankton support the zooplankton population. Which statement resolves this apparent contradiction?} \hfill [Tag: Ecological pyramids | Hard]
\begin{enumerate}[label=(\alph*),topsep=2pt,itemsep=1pt,partopsep=0pt,parsep=0pt]
    \item {The energy pyramid is inverted because zooplankton create energy during assimilation}
    \item {The productivity pyramid is inverted because consumers reproduce faster than producers}
    \item {The number pyramid must be upright because each phytoplankton cell is larger than a grazer}
    \item {The biomass pyramid is inverted because phytoplankton have rapid turnover and high productivity}
\end{enumerate}
\vspace{2mm}

\noindent\textbf{Q59. A grassland has producer net production of 12 , 000 kJ m − 2 y r − 1 . Transfer efficiencies are 15\% from producers to herbivores, 10\% from herbivores to primary carnivores, and 20\% from primary to secondary carnivores. How much production reaches the secondary carnivores?} \hfill [Tag: Ecological pyramids | Hard]
\begin{enumerate}[label=(\alph*),topsep=2pt,itemsep=1pt,partopsep=0pt,parsep=0pt]
    \item {18 kJ m − 2}
    \item {360 kJ m − 2}
    \item {180 kJ m − 2}
    \item {36 kJ m − 2}
\end{enumerate}
\vspace{2mm}

\noindent\textbf{Q60. Two sites are abandoned simultaneously: Site X is exposed bare rock after glacial retreat, while Site Y is an abandoned cultivated field with intact soil and a surviving seed bank. Which comparison is most accurate?} \hfill [Tag: Ecological succession | Hard]
\begin{enumerate}[label=(\alph*),topsep=2pt,itemsep=1pt,partopsep=0pt,parsep=0pt]
    \item {Both sites undergo secondary succession because each previously supported living organisms}
    \item {Site Y undergoes primary succession because cultivation removed the former vegetation}
    \item {Site X undergoes secondary succession and usually develops faster than Site Y}
    \item {Site X undergoes primary succession and usually develops more slowly than Site Y}
\end{enumerate}
\vspace{2mm}


\section{Section B: Comprehensive Subjective Problem with Detailed Analytical Rubric}

\begin{subjectiveq}{2}{1}
\textbf{Trophic Energetics, Ecological Pyramids and Biomagnification:}\\A limnological study of an open freshwater lake ecosystem determines that incident solar radiation over a surface area of $10{,}000\text{ m}^2$ is $6.0 \times 10^9\text{ kcal/year}$. The phytoplankton photosynthetic community has a Gross Primary Productivity (GPP) efficiency of $1.5\%$. Autotrophic respiration ($R_a$) consumes $40\%$ of the fixed gross chemical energy.\\\\(a) Calculate the Gross Primary Productivity (GPP) and Net Primary Productivity (NPP) in $\text{kcal/year}$ and $\text{kcal/m}^2/\text{year}$.\\(b) Assuming Lindeman's 10\% thermodynamic ecological transfer efficiency applies strictly between consecutive trophic levels, compute the energy available to: (i) Herbivorous Zooplankton (Primary Consumers), (ii) Small Carnivorous Planktivorous Fish (Secondary Consumers), and (iii) Apex Piscivorous Fish/Osprey (Tertiary Consumers).\\(c) If a persistent organochlorine pesticide (DDT) is present in the lake water at an ambient concentration of $0.00005\text{ ppm}$ ($0.05\text{ ppb}$) and undergoes bioconcentration by a factor of $800$ in phytoplankton, with a trophic biomagnification factor of $5\times$ at each subsequent consumer step, calculate the DDT concentration in the apex predator. State why the inverted pyramid of biomass occurs in open oceans but the pyramid of energy is always upright.
\end{subjectiveq}

\begin{solbox}
\textbf{Part (a): Solar Energy Partitioning and Net Primary Production:}\\Given total incident solar insolation $E_{\text{solar}} = 6.0 \times 10^9\text{ kcal/year}$ across $A = 10{,}000\text{ m}^2$:\begin{align*}\text{GPP} &= E_{\text{solar}} \times 0.015 = 6.0 \times 10^9 \times 0.015 = 9.0 \times 10^7\text{ kcal/year} \\\text{GPP per unit area} &= \frac{9.0 \times 10^7}{10{,}000} = 9{,}000\text{ kcal/m}^2/\text{year}\end{align*}Autotrophic respiration consumes $40\%$ of GPP ($R_a = 0.40 \times \text{GPP}$):\begin{align*}\text{NPP} &= \text{GPP} - R_a = \text{GPP} \times (1 - 0.40) = 0.60 \times 9.0 \times 10^7 = 5.4 \times 10^7\text{ kcal/year} \\\text{NPP per unit area} &= \frac{5.4 \times 10^7}{10{,}000} = 5{,}400\text{ kcal/m}^2/\text{year}.\end{align*}\textbf{Part (b): Lindeman's Trophic Efficiency Energy Flow:}\\Applying Lindeman's $10\%$ rule ($\lambda = 0.10$):\begin{align*}E_1\text{ (Producers/NPP)} &= 5.4 \times 10^7\text{ kcal/year} \\E_2\text{ (Primary Consumers - Zooplankton)} &= E_1 \times 0.10 = 5.4 \times 10^6\text{ kcal/year} \\E_3\text{ (Secondary Consumers - Small Fish)} &= E_2 \times 0.10 = 5.4 \times 10^5\text{ kcal/year} \\E_4\text{ (Tertiary Consumers - Apex Piscivores)} &= E_3 \times 0.10 = 5.4 \times 10^4\text{ kcal/year}.\end{align*}\textbf{Part (c): Biomagnification Calculation and Inverted Pyramids:}\\Bioconcentration into phytoplankton (Producers):\begin{equation*}C_{\text{phytoplankton}} = C_{\text{water}} \times 800 = 0.00005\text{ ppm} \times 800 = 0.04\text{ ppm}.\end{equation*}Applying subsequent biomagnification factors ($BMF = 5$ at each step):\begin{align*}C_{\text{zooplankton}} &= 0.04\text{ ppm} \times 5 = 0.20\text{ ppm} \\C_{\text{small fish}} &= 0.20\text{ ppm} \times 5 = 1.00\text{ ppm} \\C_{\text{apex predator}} &= 1.00\text{ ppm} \times 5 = 5.00\text{ ppm}.\end{align*}The concentration increases by a factor of $100{,}000$ from ambient water to the apex predator.\\\textit{Pyramid Inversion Principle:} In marine ecosystems, phytoplankton possess an exceptionally rapid turnover rate (days) and high productivity but low standing biomass, whereas long-lived zooplankton and baleen whales have slower turnover, creating an \textbf{inverted pyramid of biomass}. In contrast, the \textbf{pyramid of energy is ALWAYS upright} in all ecosystems because the Second Law of Thermodynamics dictates irreversible entropy generation and metabolic heat dissipation at each consumer transfer step.
\end{solbox}

\begin{examinertip}
\textbf{Marking Scheme Breakdown (Total: 10 Marks):}\\
\textbullet\ \textbf{3 Marks:} Clear formulation of governing environmental identities, equations, and thermodynamic/constitutional foundations.\\
\textbullet\ \textbf{4 Marks:} Rigorous algebraic/numerical substitution, dimensional consistency, and unit handling.\\
\textbullet\ \textbf{3 Marks:} Critical evaluation of environmental implications, regulatory criteria (CPCB/UNDRR/IUCN), and policy mitigations.
\end{examinertip}
